\chapter*{Use Cases} \label{use_cases}

\section{AI Accelerator Weight Memory Protection}

Deep neural network inference accelerators store large weight matrices in high-density on-chip SRAM or off-chip DRAM. These memories are increasingly operated at reduced supply voltage to minimize power consumption, raising the BER above the correction capacity of conventional on-chip ECC (typically rated for 1--8 errors). A single bit-flip in the exponent of a FP16 weight can change its value by a factor of $2^{1024}$, causing immediate and catastrophic accuracy collapse in the affected network layer.

[ACRONYM] is particularly well suited to this application for three reasons. First, DNN inference is tolerant of small residual errors---model accuracy degrades gracefully with occasional uncorrected flips---so the $\geq$95\% probabilistic correction guarantee of [ACRONYM] is sufficient. Second, [ACRONYM]'s 100\% overhead at 4,096-bit blocks (one metadata SRAM bank per data bank) is substantially lower than BCH's 174\% at equivalent BER, reducing area and access bandwidth. Third, weight memories are read-dominated (weights are loaded once per inference pass), so the correction latency is amortized over many inference operations.

In this deployment, each weight tile of $L$ bits (e.g., a 4,096-bit row of a layer's weight matrix) is encoded by [ACRONYM] before being written to the weight SRAM. During inference, each tile is read, decoded, and corrected before being presented to the multiply-accumulate (MAC) units. The ECC engine operates inline with the weight SRAM controller, adding one correction latency per tile read.

\section{NAND Flash Storage}

NAND flash memory exhibits increasing BER as cells age through program-erase cycles. MLC (2 bits/cell) flash begins life at BERs below $10^{-5}$ and reaches $10^{-2}$ near end-of-life; TLC (3 bits/cell) flash reaches $10^{-1}$ at end-of-life in some configurations. Additionally, flash page-disturb operations (when a page in the same block is programmed) cause burst errors: contiguous cells within a page can flip together, producing runs of 10--100 consecutive flipped bits that overwhelm BCH codes rated for random errors.

[ACRONYM] addresses both failure modes simultaneously. The Feistel permutation scatters burst errors uniformly, and the constraint-guided decoder corrects the resulting distributed errors with the same efficiency as random errors. The $O(1/\sqrt{L})$ overhead scaling is particularly attractive as flash page sizes increase from 4 KB to 16 KB and beyond: at 16 KB pages with CRC-32, [ACRONYM] metadata overhead drops to $\sim$50\%, compared to BCH which maintains roughly constant overhead per bit regardless of page size.

In this deployment, [ACRONYM] operates as the outer ECC layer in a multi-level ECC architecture, with a lightweight inner code (e.g., a simple Hamming code) handling low-BER early-life operation, and [ACRONYM] activated when the inner code reports correction failures beyond its capacity.

\section{DRAM Reliability at Reduced Supply Voltage}

DRAM cells exhibit exponentially increasing soft error rates (from cosmic-ray-induced charge deposition) and hard error rates (from retention failures) as supply voltage is reduced below nominal for energy savings. Memory systems in data centers and mobile devices increasingly employ aggressive voltage scaling, pushing BERs into ranges where conventional on-chip ECC (typically $t=1$ single-error-correcting codes) is insufficient.

[ACRONYM] provides a drop-in enhancement to the DRAM ECC engine. For a standard 64-byte DRAM burst (512 bits), [ACRONYM] with CRC-16 (50\% overhead, 256 metadata bits) corrects up to $\sim$25 bit-flips, covering the expected error count at aggressively scaled voltages with comfortable margin. The hardware decoder (parallel CRC units + combinatorial search) adds minimal area to the DRAM controller and operates within the memory read latency budget.

In this deployment, the 512-bit data burst and 256-bit [ACRONYM] metadata are stored in adjacent DRAM addresses. On a read, the data and metadata are fetched simultaneously, and the [ACRONYM] decoder produces the corrected burst within the burst transfer window.

\section{General-Purpose Data Integrity for High-BER Channels}

Beyond specific memory technologies, [ACRONYM] applies to any digital storage or communication channel where: (i) the BER exceeds the correction capacity of conventional lightweight ECC, (ii) burst errors are present, (iii) Galois field arithmetic is not available or desirable, or (iv) the overhead of BCH codes is prohibitive. Representative additional applications include:

\begin{itemize}
    \item \textbf{Optical storage:} CD/DVD/Blu-ray discs with physical damage produce burst errors of tens to hundreds of bits; [ACRONYM]'s Feistel scattering handles these natively.
    \item \textbf{Magnetic HDD:} Sector-level ECC on spinning disks at high areal densities encounters BERs in the $10^{-4}$--$10^{-3}$ range in severe head-disk interface conditions.
    \item \textbf{Radiation-hardened computing:} Satellite and aerospace electronics exposed to high-energy particle strikes require ECC capable of handling multi-bit upset events; [ACRONYM]'s scalable correction capability at tunable overhead directly addresses this need.
    \item \textbf{In-memory computing:} Compute-in-memory (CIM) architectures that perform arithmetic directly in SRAM or ReRAM arrays introduce systematic bit-flip noise from analog imprecision; [ACRONYM] can protect the operand values stored in these arrays.
\end{itemize}
